
\documentclass[template=tabling,81pt,headonall]{azmoon}
\usepackage{xepersian}
\usepackage{amsfonts}
\usepackage{graphicx}
\usepackage{svg}
\svgpath{ {./images/} }
\graphicspath{ {./images/} }
\settextfont{Yas}
\setdigitfont{A Iranian Sans}
\usepackage{fontawesome5}

\printanswers
    \teacher{محمد صالح علی اکبری}
    \teachertitle{دبیر}
    \city{گناباد}
    \schooltitle{دبیرستان دوره اول}
    \school{شهید جوانبخت}
    \grade{ریاضی هفتم (نمونه ۱)}
    \branch{نمونه ۱}
    \topic{فصل ۱: راهبردهای حل مسئله}
    \examdate{۲۳ مهر ۱۴۰۵}
    \answertime{۴۰ دقیقه}
    \begin{document}
	\begin{questions}
		\nointerlineskip%
		\vskip-\baselineskip
		\question{%
در الگوی زیر، چوب‌کبریت‌ها برای ساختن مربع‌های به هم چسبیده به کار رفته‌اند:

\noindent\textbf{شکل (۱):} $4$ چوب‌کبریت \hfill \textbf{شکل (۲):} $7$ چوب‌کبریت \hfill \textbf{شکل (۳):} $10$ چوب‌کبریت

الف) شکل چهارم با چند چوب‌کبریت ساخته می‌شود؟ \\
ب) رابطهٔ جبری بین شمارهٔ شکل ($n$) و تعداد چوب‌کبریت‌ها را بنویسید. \\
ج) شکل بیستم ($20$) از چند چوب‌کبریت تشکیل شده است؟‌
\\‌
\\‌
\\}\question{%
کشاورزی نصف زمین خود را گندم کاشته و $\frac{1}{3}$ از باقی‌ماندهٔ زمین را جو کاشته است. اگر پس از این کارها، مساحت زمین باقی‌مانده برابر با $400$ متر مربع باشد، کل مساحت زمین کشاورزی چقدر است؟ \textbf{(حل مسئله با «رسم شکل» و نوشتن توضیحات الزامی است)}‌
\\‌
\\‌
\\‌
\\‌
\\}\question{%
در یک پارکینگ در مجموع $18$ دستگاه خودرو و موتورسیکلت پارک شده است. اگر تعداد کل چرخ‌های آن‌ها روی هم $56$ عدد باشد، در این پارکینگ چند خودرو و چند موتورسیکلت وجود دارد؟ \textbf{(جدول حدس و آزمایش را با حداقل سه حدس رسم کرده و پاسخ نهایی را بنویسید)}‌
\\‌
\\‌
\\‌
\\‌
\\}\question{%
حاصل‌ضرب سن دو برادر برابر $48$ است. اگر بدانیم مجموع سن آن‌ها کمترین مقدار ممکن است، سن هر کدام از برادرها چند سال است؟ \textbf{(جدول حالت‌های مختلف ضرب را به صورت نظام‌دار بنویسید و حالت مطلوب را مشخص کنید)}‌
\\‌
\\‌
\\‌
\\}\end{questions}
    \end{document}
    